\textbf{Phần 1. Trắc nghiệm lựa chọn}
\setcounterpageref{ex}{0}

\begin{ex} %[2D14V72]
Hình tam giác có độ dài ba cạnh lần lượt là $15\text{ cm}, 20\text{ cm}, 25\text{ cm}$. Chu vi hình tam giác đó là
\choice
{$50\text{ cm}$}
{$55\text{ cm}$}
{\True $60\text{ cm}$}
{$65\text{ cm}$}
\loigiai{
Chu vi: $15 + 20 + 25 = 60$ cm.
\textbf{Đáp án của câu hỏi là C}
}
\end{ex}

\begin{ex} %[2D14V72]
Hình tứ giác có độ dài bốn cạnh đều bằng $5\text{ cm}$. Chu vi của hình tứ giác đó là
\choice
{$15\text{ cm}$}
{\True $20\text{ cm}$}
{$25\text{ cm}$}
{$10\text{ cm}$}
\loigiai{
Chu vi: $5 \times 4 = 20$ cm.
\textbf{Đáp án của câu hỏi là B}
}
\end{ex}

\begin{ex} %[2D14V72]
Đường gấp khúc gồm $3$ đoạn thẳng có độ dài lần lượt là $12\text{ cm}, 18\text{ cm}, 30\text{ cm}$. Độ dài đường gấp khúc đó bằng bao nhiêu đề-xi-mét?
\choice
{$5\text{ dm}$}
{$60\text{ dm}$}
{$7\text{ dm}$}
{\True $6\text{ dm}$}
\loigiai{
Độ dài: $12 + 18 + 30 = 60\text{ cm} = 6\text{ dm}$.
\textbf{Đáp án của câu hỏi là D}
}
\end{ex}

\begin{ex} %[2D14V72]
Hình nào dưới đây có thể lăn được về mọi hướng?
\choice
{\True Khối cầu}
{Khối trụ}
{Khối lập phương}
{Khối hộp chữ nhật}
\loigiai{
Khối cầu có bề mặt cong hoàn toàn nên lăn được về mọi hướng.
\textbf{Đáp án của câu hỏi là A}
}
\end{ex}

\begin{ex} %[2D14V72]
Khối trụ có hai mặt đáy là hình gì?
\choice
{Hình vuông}
{Hình tam giác}
{\True Hình tròn}
{Hình chữ nhật}
\loigiai{
Hai mặt đáy của khối trụ là hai hình tròn bằng nhau.
\textbf{Đáp án của câu hỏi là C}
}
\end{ex}

\begin{ex} %[2D14V72]
Cho ba điểm $A, B, C$ thẳng hàng, điểm $B$ nằm giữa $A$ và $C$. Biết $AB = 8\text{ cm}, BC = 12\text{ cm}$. Độ dài đoạn thẳng $AC$ là
\choice
{$4\text{ cm}$}
{\True $20\text{ cm}$}
{$16\text{ cm}$}
{$22\text{ cm}$}
\loigiai{
$AC = AB + BC = 8 + 12 = 20$ cm.
\textbf{Đáp án của câu hỏi là B}
}
\end{ex}

\begin{ex} %[2D14V72]
Một hình chữ nhật có chiều dài $16\text{ cm}$, chiều rộng ngắn hơn chiều dài $6\text{ cm}$. Nửa chu vi hình chữ nhật đó là
\choice
{$22\text{ cm}$}
{$52\text{ cm}$}
{$20\text{ cm}$}
{\True $26\text{ cm}$}
\loigiai{
Chiều rộng: $16 - 6 = 10$ cm. Nửa chu vi: $16 + 10 = 26$ cm.
\textbf{Đáp án của câu hỏi là D}
}
\end{ex}

\begin{ex} %[2D14V72]
Một đoạn dây đồng dài $32\text{ cm}$ được uốn thành một hình tứ giác có bốn cạnh bằng nhau. Mỗi cạnh của hình tứ giác dài bao nhiêu xăng-ti-mét?
\choice
{\True $8\text{ cm}$}
{$6\text{ cm}$}
{$7\text{ cm}$}
{$9\text{ cm}$}
\loigiai{
$32 : 4 = 8$ cm.
\textbf{Đáp án của câu hỏi là A}
}
\end{ex}

\begin{ex} %[2D14V72]
Khối hộp chữ nhật có tất cả bao nhiêu mặt?
\choice
{$4$ mặt}
{$8$ mặt}
{\True $6$ mặt}
{$12$ mặt}
\loigiai{
Khối hộp chữ nhật có 6 mặt.
\textbf{Đáp án của câu hỏi là C}
}
\end{ex}

\begin{ex} %[2D14V72]
Một hình tam giác có chu vi $45\text{ cm}$, hai cạnh dài $15\text{ cm}$ và $18\text{ cm}$. Độ dài cạnh còn lại là
\choice
{$10\text{ cm}$}
{\True $12\text{ cm}$}
{$14\text{ cm}$}
{$13\text{ cm}$}
\loigiai{
Cạnh thứ ba: $45 - (15 + 18) = 12$ cm.
\textbf{Đáp án của câu hỏi là B}
}
\end{ex}

\begin{ex} %[2D14V72]
Vật nào sau đây có dạng khối trụ?
\choice
{Hộp diêm}
{Quả bóng đá}
{Viên gạch}
{\True Lon sữa bò}
\loigiai{
Lon sữa bò có dạng khối trụ.
\textbf{Đáp án của câu hỏi là D}
}
\end{ex}

\begin{ex} %[2D14V72]
Cho $4$ điểm phân biệt trong đó không có ba điểm nào thẳng hàng. Nối từng cặp điểm thì vẽ được bao nhiêu đoạn thẳng?
\choice
{\True $6$ đoạn thẳng}
{$4$ đoạn thẳng}
{$5$ đoạn thẳng}
{$8$ đoạn thẳng}
\loigiai{
Số đoạn thẳng: $4 \times 3 : 2 = 6$ đoạn thẳng.
\textbf{Đáp án của câu hỏi là A}
}
\end{ex}

\textbf{Phần 2. Trắc nghiệm đúng sai}
\setcounterpageref{ex}{0}

\begin{ex} %[2D14V72]
Xét tính đúng sai của các khẳng định về hình học:
\choiceTF
{\True Đường thẳng không bị giới hạn ở hai đầu.}
{Hình tứ giác là hình có 3 cạnh và 3 đỉnh.}
{\True Chu vi của một hình bằng tổng độ dài các cạnh bao quanh hình đó.}
{\True Khối trụ có thể lăn được khi đặt nằm ngang.}
\loigiai{
\begin{itemize}
    \item Ý a) Đúng: Đúng định nghĩa đường thẳng.
    \item Ý b) Sai: Hình tứ giác có 4 cạnh và 4 đỉnh. Sai.
    \item Ý c) Đúng: Đúng quy tắc tính chu vi.
    \item Ý d) Đúng: Đúng đặc điểm khối trụ.
\end{itemize}
\begin{itemize}
    \item a) \textbf{Đúng}.
    \item b) \textbf{Sai}.
    \item c) \textbf{Đúng}.
    \item d) \textbf{Đúng}.
\end{itemize}
}
\end{ex}

\begin{ex} %[2D14V72]
Một hình tam giác có độ dài các cạnh là 12 cm, 16 cm và 20 cm:
\choiceTF
{\True Chu vi của hình tam giác là 48 cm.}
{\True Cạnh dài nhất dài hơn cạnh ngắn nhất 8 cm.}
{Chu vi của tam giác lớn hơn 5 dm.}
{\True Nếu uốn một sợi dây thành tam giác này thì cần sợi dây dài ít nhất 48 cm.}
\loigiai{
\begin{itemize}
    \item Ý a) Đúng: 12 + 16 + 20 = 48 cm. Đúng.
    \item Ý b) Đúng: 20 - 12 = 8 cm. Đúng.
    \item Ý c) Sai: 48 cm < 50 cm = 5 dm. Sai.
    \item Ý d) Đúng: Đúng.
\end{itemize}
\begin{itemize}
    \item a) \textbf{Đúng}.
    \item b) \textbf{Đúng}.
    \item c) \textbf{Sai}.
    \item d) \textbf{Đúng}.
\end{itemize}
}
\end{ex}

\textbf{Phần 3. Trả lời ngắn}
\setcounterpageref{ex}{0}

\begin{ex} %[2D14V72]
Tính chu vi hình tam giác có độ dài ba cạnh đều bằng 9 cm.
\shortans{27}
\loigiai{
\textbf{Đáp số:} $27$\\
$9 \times 3 = 27$ cm.
}
\end{ex}

\begin{ex} %[2D14V72]
Một hình tứ giác có bốn cạnh dài lần lượt là 10 cm, 12 cm, 15 cm, 13 cm. Chu vi của tứ giác là bao nhiêu đề-xi-mét?
\shortans{5}
\loigiai{
\textbf{Đáp số:} $5$\\
$10 + 12 + 15 + 13 = 50\text{ cm} = 5$ dm.
}
\end{ex}

\begin{ex} %[2D14V72]
Một đường gấp khúc gồm 4 đoạn thẳng bằng nhau, mỗi đoạn dài 5 cm. Độ dài đường gấp khúc là bao nhiêu xăng-ti-mét?
\shortans{20}
\loigiai{
\textbf{Đáp số:} $20$\\
$5 \times 4 = 20$ cm.
}
\end{ex}

\begin{ex} %[2D14V72]
Khối lập phương có bao nhiêu đỉnh?
\shortans{8}
\loigiai{
\textbf{Đáp số:} $8$\\
Khối lập phương có 8 đỉnh.
}
\end{ex}

\begin{ex} %[2D14V72]
Một mảnh đất hình chữ nhật có chiều dài 25 m, chiều rộng 15 m. Chu vi mảnh đất là bao nhiêu mét?
\shortans{80}
\loigiai{
\textbf{Đáp số:} $80$\\
$(25 + 15) \times 2 = 80$ m.
}
\end{ex}

\begin{ex} %[2D14V72]
Đoạn thẳng AB dài 18 cm, điểm M là điểm chia đôi đoạn AB thành 2 đoạn bằng nhau. Độ dài đoạn AM là bao nhiêu xăng-ti-mét?
\shortans{9}
\loigiai{
\textbf{Đáp số:} $9$\\
$18 : 2 = 9$ cm.
}
\end{ex}
